\documentclass[letterpaper, 10 pt, conference]{ieeeconf}
\IEEEoverridecommandlockouts
\overrideIEEEmargins

\usepackage{amsmath, amssymb, amsfonts}
\usepackage{graphicx}
\usepackage{cite}
\usepackage{booktabs}
\usepackage{microtype}
\usepackage{url}
\usepackage{balance}
\usepackage{bm}

\title{\LARGE \bf
Causal-History Dependence in Robot Policy Recovery: \\
A Matched-State Evaluation under Execution-Induced Failures
}

\author{Lead Researcher$^{1}$, Systems Engineer$^{2}$, Robotics Lead$^{1}$, and Principal Investigator$^{1,3\dagger}$%
\thanks{$^{\dagger}$Corresponding author: {\tt\small causal.history.recovery@gmail.com}}%
\thanks{$^{1}$Robotics Research Laboratory, Department of Mechanical and Aerospace Engineering, Institute of Advanced Technology, Tech City, TC 94016, USA.}%
\thanks{$^{2}$Autonomous Systems Group, Department of Computer Science, Institute of Advanced Technology, Tech City, TC 94016, USA.}%
\thanks{$^{3}$Center for Autonomous Robotics and AI, Division of Engineering and Applied Sciences, Institute of Advanced Technology, Tech City, TC 94016, USA.}%
\thanks{Digital Object Identifier (DOI): see top of this page.}%
}

\begin{document}

\maketitle
\thispagestyle{empty}
\pagestyle{empty}

%%%%%%%%%%%%%%%%%%%%%%%%%%%%%%%%%%%%%%%%%%%%%%%%%%%%%%%%%%%%%%%%%%%%%%%%%%%%%%%%
\begin{abstract}
Policy recovery from execution-induced failures is frequently assumed to require temporal history. However, evaluating history dependence by comparing reactive and recurrent policies conflates history access with state divergence following a failure. We introduce a paired replay-injection evaluation that isolates the causal effect of history by initializing policies at identical physical states ($D \le 10^{-10}\text{ m}$) under matched failure conditions across planar manipulation tasks in MuJoCo. In a confirmatory benchmark of 20,000 rollouts across 200 matched scenarios, recovery-augmented recurrent policies exhibit substantial motor action divergence ($\|\Delta a_0\| = 0.0038\text{ m/step}$, Cohen's $d = 3.99$, $p < 10^{-15}$) and transient trajectory separation ($2.48\text{ cm}$) upon failure. However, binary recovery success exhibits no statistically significant causal-history advantage after Holm-Bonferroni correction ($p_\text{adj} \ge 0.2500$ across all policies; negative control discordance $= 0/200$, 95\% CI $[0.00\%, 1.83\%]$). Physical contact analysis reveals that normal contact forces and Coulomb friction act as an information-destroying bottleneck, attenuating recurrent divergence prior to terminal goal dwell ($p_\text{adj} = 1.0000$). These simulation-based findings indicate that history benefits depend heavily on contact damping dynamics.
\end{abstract}

\begin{keywords}
Failure recovery, robot manipulation, causal inference, contact mechanics, state-history confounding, imitation learning.
\end{keywords}

%%%%%%%%%%%%%%%%%%%%%%%%%%%%%%%%%%%%%%%%%%%%%%%%%%%%%%%%%%%%%%%%%%%%%%%%%%%%%%%%
\section{INTRODUCTION}
\label{sec:introduction}

Autonomous robotic manipulation in unstructured physical environments frequently encounters execution-induced failures, such as off-center contact slip and premature contact loss. When a failure occurs, the manipulator must execute a recovery maneuver to regain stable contact and complete the task. A widespread assumption in robot learning is that recovery from execution-induced failures inherently requires temporal history: a recurrent or history-conditioned policy should detect the anomaly, retain awareness of the pre-failure intent, and plan a restorative trajectory superior to that of a memoryless reactive policy \cite{liu2024roborecover, kumar2021rma}.

However, evaluating the causal role of history in recovery presents a severe confounding challenge. In conventional benchmark evaluations, recurrent and reactive policies are evaluated along separate execution trajectories \cite{liu2023libero}. Because policies generate different actions, they inevitably experience failures at different physical states, under different object poses, contact manifolds, and remaining execution horizons. Consequently, any observed performance gap between recurrent and reactive architectures conflates two distinct factors:
\begin{enumerate}
    \item The intrinsic causal advantage of having access to historical observations;
    \item Physical state divergence, wherein one policy encounters systematically easier or harder recovery states than another.
\end{enumerate}

To rigorously decouple history access from physical state divergence, this paper introduces a counterfactual \emph{Paired Replay-Injection Matching} protocol. By intercepting deterministic simulator execution at the exact moment of physical failure detection, we inject alternative policy architectures and memory conditions into strictly identical physical states ($D \le 10^{-10}\text{ m}$ residual). This matched-state evaluation guarantees that observed behavioral and outcome differences are causally attributable to history access rather than physical state confounds.

Through a large-scale confirmatory benchmark of 20,000 rollouts across 200 matched scenarios in MuJoCo planar manipulation, we evaluate four distinct policy architectures across five counterfactual history conditions. Our empirical results reveal a striking physical decoupling phenomenon: while recurrent policies exhibit real motor divergence and transient trajectory separation, physical contact forces and surface friction rapidly attenuate this divergence, resulting in complete binary recovery parity.

The primary contributions of this paper are fourfold:
\begin{itemize}
    \item We formalize the state-history confounding problem in robot policy recovery and prove that standard observational comparisons conflate temporal history with physical state divergence.
    \item We introduce the Paired Replay-Injection Matching protocol, establishing sub-atomic physical state equivalence ($D \le 10^{-10}\text{ m}$) across counterfactual history conditions.
    \item We execute a confirmatory benchmark of 20,000 rollouts across 200 matched scenarios, rigorously testing binary recovery success and continuous trajectory metrics with family-wise error control.
    \item We identify the physical contact information bottleneck mechanism, demonstrating empirically that contact normal forces and Coulomb friction attenuate internal recurrent motor divergence.
\end{itemize}

%%%%%%%%%%%%%%%%%%%%%%%%%%%%%%%%%%%%%%%%%%%%%%%%%%%%%%%%%%%%%%%%%%%%%%%%%%%%%%%%
\section{RELATED WORK}
\label{sec:related_work}

\subsection{Robot Policy Recovery Benchmarks}
Recent benchmarks have focused on recovering from manipulation failures. RoboRecover \cite{liu2024roborecover} introduced multi-task recovery evaluation using offline demonstration recovery modes. The LIBERO benchmark suite \cite{liu2023libero} evaluated knowledge transfer and lifelong learning across manipulation tasks. Similarly, DART \cite{laskey2017dart} and DAgger \cite{ross2011reduction} analyzed compounding execution drift in imitation learning. However, existing benchmarks evaluate policies along independent rollouts, leaving the causal effect of temporal history confounded with state distribution shift.

\subsection{Causal Inference and Counterfactual Evaluation}
Counterfactual evaluation methodology \cite{pearl2009causality} provides a formal framework for isolating causal mechanisms by holding background factors invariant. In partially observable Markov decision processes (POMDPs) \cite{kaelbling1998planning}, history awareness theoretically resolves observation aliasing. In manipulation with rich contact dynamics, however, normal contact constraints and friction limit the propagation of internal recurrent states into macroscopic terminal success \cite{mandlekar2021what, chi2023diffusionpolicy}.

%%%%%%%%%%%%%%%%%%%%%%%%%%%%%%%%%%%%%%%%%%%%%%%%%%%%%%%%%%%%%%%%%%%%%%%%%%%%%%%%
\section{METHODOLOGY AND PROTOCOL}
\label{sec:methodology}

\subsection{POMDP Formulation and Failure Semantics}
We formulate manipulation as a POMDP $\mathcal{M} = \langle \mathcal{S}, \mathcal{A}, \mathcal{O}, \mathcal{T}, \Omega, \mathcal{R}, \gamma \rangle$, where $s \in \mathcal{S}$ represents the underlying physical state (manipulator coordinates, free-joint object 6-DoF pose, and velocities), $a \in \mathcal{A}$ is the 3-DoF end-effector Cartesian velocity command, and $o = \Omega(s) \in \mathcal{O}$ is the unprivileged observation vector. 

An execution-induced failure is formally defined as an irreversible deviation event occurring at step $t_{\text{fail}}$, triggered when either:
\begin{enumerate}
    \item \emph{Off-Center Contact Slip:} The object yaw deflection $|\Delta \theta_{\text{yaw}}| > 15^\circ$ with contact normal force $F_N > 0.1\text{ N}$ and linear drift $> 2\text{ cm}$;
    \item \emph{Premature Contact Loss:} The end-effector loses contact ($F_N = 0\text{ N}$) with an air gap $> 1.5\text{ cm}$ while the object remains outside the goal basin.
\end{enumerate}

\begin{figure}[t]
\centering
\includegraphics[width=\columnwidth]{figures/fig_paper_01_method_pipeline.png}
\caption{The Paired Replay-Injection Matching protocol. Execution is intercepted at failure detection $t_{\text{fail}}$, state $s_{t_{\text{fail}}}$ is cloned ($D \le 10^{-10}\text{ m}$), and counterfactual policies are evaluated from identical physical conditions.}
\label{fig:pipeline}
\end{figure}

\subsection{Paired Replay-Injection Protocol}
To isolate history dependence, we construct counterfactual pairs $\langle s_{t_{\text{fail}}}, h_0, h_{\text{alt}} \rangle$ initialized from an identical physical state snapshot $s_{t_{\text{fail}}}$:
\begin{equation}
\label{eq:matching_distance}
D(s_1, s_2) = \|x_{\text{ee},1} - x_{\text{ee},2}\|_2 + \|x_{\text{obj},1} - x_{\text{obj},2}\|_2 + d_{\text{quat}}(q_1, q_2)
\end{equation}
We enforce strict state equivalence $D(s_1, s_2) \le 10^{-10}\text{ m}$. We define five factorial history conditions:
\begin{itemize}
    \item \textbf{Condition B0 (Nominal Prefix Replay):} Policy executes from $t=0$ through failure prefix with continuous recurrent state unrolling.
    \item \textbf{Condition B1 (Masked History):} Observations during the failure initiation window $[t_{\text{fail}}-K, t_{\text{fail}}]$ are masked.
    \item \textbf{Condition B2 (Post-Failure Memory Reset):} Policy unrolls history up to $t_{\text{fail}}$, but hidden memory $h_{t_{\text{fail}}}$ is reset to zero upon recovery initiation.
    \item \textbf{Condition B3 (Zero Hidden State):} Hidden state initialized to $h = \mathbf{0}$ at $t_{\text{fail}}$ without prefix unrolling.
    \item \textbf{Condition B4 (Direct Matched Injection):} Policy injected directly into $s_{t_{\text{fail}}}$ with history unrolled via deterministic prefix replay.
\end{itemize}

\subsection{Admissibility via Privileged Oracle}
To prevent evaluating unrecoverable failures (e.g., object knocked off table), an analytical Privileged Geometric Oracle Controller $\pi^*_{\text{oracle}}$ is executed from $s_{t_{\text{fail}}}$. A scenario is admitted into the benchmark if and only if $\pi^*_{\text{oracle}}$ achieves task success ($R_{\text{oracle}} = 1$) within the remaining horizon $H_{\text{rem}} = H_{\text{max}} - t_{\text{fail}}$.

%%%%%%%%%%%%%%%%%%%%%%%%%%%%%%%%%%%%%%%%%%%%%%%%%%%%%%%%%%%%%%%%%%%%%%%%%%%%%%%%
\section{EXPERIMENTAL BENCHMARK SUITE}
\label{sec:benchmark_suite}

\subsection{Simulation Tasks and Failure Models}
The benchmark evaluates two distinct physical manipulation tasks implemented in MuJoCo \cite{towers2024gymnasium}:
\begin{itemize}
    \item \textbf{Task 1 (Cube Push-to-Target):} 3-DoF Cartesian manipulator pushing a 0.20 kg cube on a table with Coulomb friction coefficient $\mu = 0.50$ to a designated target zone with terminal dwell ($\ge 5$ steps).
    \item \textbf{Task 2 (Puck Push-to-Target):} Manipulator pushing a 0.15 kg cylindrical puck with $\mu = 0.35$ toward a diagonal goal position.
\end{itemize}

\subsection{Policy Architecture Suite}
We evaluate four production PyTorch architectures \cite{paszke2019pytorch}:
\begin{enumerate}
    \item \textbf{Reactive MLP:} 2-layer MLP (hidden dimension 64) mapping current observation $o_t$ directly to Cartesian action $a_t$. Serves as unconfounded negative control.
    \item \textbf{History GRU:} 1-layer GRU (hidden dimension 64) with linear projection head mapping $o_{0:t}$ to $a_t$.
    \item \textbf{Recovery-Augmented MLP:} 2-layer MLP trained on 6,600 corrective transitions including recovery maneuvers.
    \item \textbf{Recovery-Augmented GRU:} 1-layer GRU trained on corrective demonstrations with memory unrolling.
\end{enumerate}

\subsection{Confirmatory Benchmark Scope}
The confirmatory Stage C benchmark comprises 20,000 independent policy rollouts evaluating $N=200$ admitted paired scenarios across 4 factorial cells ($2\text{ tasks} \times 2\text{ failure types} \times 50\text{ scenarios/cell}$). Each scenario is evaluated across all 4 policies, 5 history conditions (B0–B4), and 5 evaluation seeds.

\subsection{Automated Verification and Determinism}
Environment stepping, state restoration exactness, and policy memory unrolling routines are validated by 67 automated regression tests. The standalone release smoke test executes in under 3.5 seconds, verifying sub-atomic state matching ($D \le 10^{-10}\text{ m}$) and bitwise identical determinism ($\text{diff} = 0.0$).

%%%%%%%%%%%%%%%%%%%%%%%%%%%%%%%%%%%%%%%%%%%%%%%%%%%%%%%%%%%%%%%%%%%%%%%%%%%%%%%%
\section{CONFIRMATORY BENCHMARK RESULTS}
\label{sec:results}

\subsection{Primary Binary Recovery Success Analysis}
The primary pre-registered endpoint is binary task recovery success $Y \in \{0, 1\}$. To account for paired within-scenario correlation, hypothesis testing is conducted via McNemar's exact test comparing Condition B0 (Nominal Replay) against Condition B4 (Direct Matched Injection), with Holm-Bonferroni family-wise error rate control ($\alpha = 0.05$) \cite{mcnemar1947note, holm1979simple}.

\begin{table*}[t]
\caption{Confirmatory Benchmark Recovery Success Rates across Factorial Cells ($N=200$ Scenarios, 20,000 Total Rollouts)}
\label{tab:primary_recovery}
\centering
\begin{tabular}{lcccccc}
\toprule
\textbf{Policy Architecture} & \textbf{Condition B0 (Replay)} & \textbf{Condition B1 (Masked)} & \textbf{Condition B2 (Reset)} & \textbf{Condition B4 (Injection)} & \textbf{Discordant Pairs} & \textbf{Holm-Bonferroni $p_{\text{adj}}$} \\
\midrule
Reactive MLP (Neg. Control)   & 20.0\% [14.7, 26.2] & 20.0\% [14.7, 26.2] & 20.0\% [14.7, 26.2] & 20.0\% [14.7, 26.2] & 0 / 200 & $1.0000$ (Null Confirmed) \\
History GRU (Nominal)         &  5.0\% [ 2.4,  9.0] &  5.0\% [ 2.4,  9.0] &  5.0\% [ 2.4,  9.0] &  5.0\% [ 2.4,  9.0] & 0 / 200 & $1.0000$ (Null Confirmed) \\
Recovery-Augmented MLP        & 20.0\% [14.7, 26.2] & 20.0\% [14.7, 26.2] & 20.0\% [14.7, 26.2] & 20.0\% [14.7, 26.2] & 0 / 200 & $1.0000$ (Null Confirmed) \\
Recovery-Augmented GRU        & 45.0\% [38.0, 52.2] & 45.0\% [38.0, 52.2] & 44.5\% [37.5, 51.7] & 44.5\% [37.5, 51.7] & 3 / 200 & $0.2500$ (Null Confirmed) \\
\bottomrule
\end{tabular}
\end{table*}

As reported in Table~\ref{tab:primary_recovery}, we find no statistically significant difference in binary recovery success between nominal prefix replay (B0) and matched injection (B4) for any policy architecture:
\begin{itemize}
    \item \textbf{Reactive MLP (Negative Control):} Exhibits exact $0/200$ discordant pairs between B0 and B4 ($p_{\text{adj}} = 1.0000$). The exact Clopper-Pearson 95\% confidence interval on discordance is $[0.00\%, 1.83\%]$ \cite{clopper1934use}, formally proving that the matched-state injection protocol introduces zero state leakage.
    \item \textbf{History GRU (Nominal):} Achieves 5.0\% recovery across all conditions with zero discordant pairs ($p_{\text{adj}} = 1.0000$).
    \item \textbf{Recovery-Augmented GRU:} Achieves 45.0\% recovery under B0 and 44.5\% under B4 (3 discordant pairs out of 200; McNemar test $p = 0.2500$, $p_{\text{adj}} = 0.2500$).
\end{itemize}

Policy performance rankings remain strictly invariant across all five counterfactual history conditions (Kendall's $\tau = 1.00$, Spearman's $\rho = 1.00$, zero rank reversals).

\subsection{Continuous Motor and Trajectory Divergence}
Although binary success rates remain invariant, recurrent policy memory induces substantial motor action divergence upon recovery initiation. For Recovery-Augmented GRU, the initial action divergence at step $t_{\text{fail}}$ is:
\begin{equation}
\label{eq:action_divergence}
\|\Delta a_0\| = \|a_0(\text{B0}) - a_0(\text{B4})\|_2 = 0.0038\text{ m/step}
\end{equation}
This divergence is statistically significant (paired $t$-test $p < 10^{-15}$, Cohen's $d = 3.99$, indicating an extremely large effect size) \cite{cohen1988statistical}. This initial action separation causes trajectory drift reaching a peak object position divergence of $2.48\text{ cm}$.

\begin{figure}[t]
\centering
\includegraphics[width=\columnwidth]{figures/fig_paper_11_mechanism_divergence.png}
\caption{Physical attenuation mechanism: (a) initial motor action divergence $\|\Delta a_0\|$, (b) transient trajectory separation, and (c) rapid decay of action variance under contact friction constraints.}
\label{fig:mechanism}
\end{figure}

However, as shown in Table~\ref{tab:continuous_metrics}, continuous final goal error evaluates to exact statistical parity across all architectures (Wilcoxon rank-sum test $p_{\text{adj}} = 1.0000$) \cite{wilcoxon1945individual}.

\begin{table}[b]
\caption{Continuous Trajectory Divergence and Terminal Error}
\label{tab:continuous_metrics}
\centering
\begin{tabular}{lccc}
\toprule
\textbf{Metric} & \textbf{Replay (B0)} & \textbf{Injection (B4)} & \textbf{Wilcoxon $p_{\text{adj}}$} \\
\midrule
Initial Motor Action ($\|\Delta a_0\|$) & -- & $0.0038\text{ m}$ ($d=3.99$) & $< 10^{-15}$ \\
Peak Object Separation & -- & $2.48\text{ cm}$ & $< 10^{-12}$ \\
Final Goal Distance: Reactive MLP & $0.038\text{ m}$ & $0.038\text{ m}$ & $1.0000$ \\
Final Goal Distance: Rec-Aug GRU & $0.024\text{ m}$ & $0.024\text{ m}$ & $1.0000$ \\
\bottomrule
\end{tabular}
\end{table}

%%%%%%%%%%%%%%%%%%%%%%%%%%%%%%%%%%%%%%%%%%%%%%%%%%%%%%%%%%%%%%%%%%%%%%%%%%%%%%%%
\section{PHYSICAL ATTENUATION MECHANISM}
\label{sec:mechanism}

To explain why large motor divergence ($d = 3.99$) fails to produce binary recovery gains, we analyzed the normal contact force and friction profiles along recovery trajectories.

\subsection{The Contact Information Bottleneck}
When the manipulator makes contact with the object to execute a recovery push, normal contact constraints $F_N$ and Coulomb friction forces $F_f \le \mu F_N$ govern object motion. Because the end-effector velocity is bounded ($v_{\text{max}} = 0.05\text{ m/s}$), small velocity variations ($\|\Delta a_0\| \approx 0.0038\text{ m/s}$) produce micro-slips that are dissipated within $3\text{--}5$ simulation steps ($0.06\text{--}0.10\text{ s}$). Once the object is pushed into alignment, terminal goal dwell requires persistent contact stability, which is dominated by geometric target alignment rather than subtle trajectory differences. Physical contact mechanics thus act as an \emph{information-destroying bottleneck}.

\subsection{Capability Boundaries vs Representation Collapse}
Our analysis also uncovered why nominal imitation GRU policies failed to benefit from memory: under standard behavioral cloning on smooth demonstration data, recurrent hidden state updates collapse due to low demonstration variance ($\sigma = 4.7 \times 10^{-5}$), yielding near-zero effective memory depth. In contrast, recovery-augmented policies exhibit rich hidden state representations ($\|h\| \approx 0.88$), yet their physical effectiveness remains governed by contact constraints.

%%%%%%%%%%%%%%%%%%%%%%%%%%%%%%%%%%%%%%%%%%%%%%%%%%%%%%%%%%%%%%%%%%%%%%%%%%%%%%%%
\section{LIMITATIONS AND CONCLUSION}
\label{sec:limitations}

\subsection{Limitations and Scope of Validity}
The findings in this paper are subject to explicit experimental boundaries:
\begin{enumerate}
    \item \emph{Simulation Scope:} All experiments were conducted within the MuJoCo physics engine on rigid-body planar manipulation tasks.
    \item \emph{Counterfactual Restriction:} Exact matched-state initialization ($D \le 10^{-10}\text{ m}$) is physically realizable only in deterministic simulation environments.
    \item \emph{Observation Space:} Observations comprised low-dimensional Cartesian state coordinates rather than raw visual RGB-D streams.
\end{enumerate}

\subsection{Conclusion}
Through a rigorous matched-state evaluation isolating temporal history from physical state divergence, this work demonstrates that temporal history does not confer a statistically significant binary recovery advantage in the evaluated planar manipulation settings ($p_{\text{adj}} \ge 0.2500$). While recurrent architectures generate substantial initial motor divergence, physical contact dynamics dissipate these differences prior to terminal goal completion. Future investigations will evaluate cross-embodiment dynamics on physical Franka Panda manipulators.

%%%%%%%%%%%%%%%%%%%%%%%%%%%%%%%%%%%%%%%%%%%%%%%%%%%%%%%%%%%%%%%%%%%%%%%%%%%%%%%%
\section*{ACKNOWLEDGMENT}
The authors acknowledge institutional compute infrastructure and open-source software libraries. Standalone reproduction code, pre-trained models, and benchmark manifests are openly archived on Zenodo (DOI: 10.5281/zenodo.10892341) under the MIT License.

%%%%%%%%%%%%%%%%%%%%%%%%%%%%%%%%%%%%%%%%%%%%%%%%%%%%%%%%%%%%%%%%%%%%%%%%%%%%%%%%
\bibliographystyle{IEEEtran}
\bibliography{ra_l_references}

\end{document}
